\subsection{Finite-support avoidance persists with real multiplication by
\texorpdfstring{$\mathcal O_5$}{O5}}
\label{sec:rm5-avoidance}

\paragraph{Finding A37 (the ordinary real-multiplication locus).}
Put $F=\mathbf Q(\sqrt5)$ and
$\mathcal O=\mathcal O_5=\mathbf Z[(1+\sqrt5)/2]$.
Consider smooth genus-two curves in characteristic two together with
a specified embedding $\mathcal O\hookrightarrow\operatorname{End}(J_C)$
fixed by the Rosati involution of the canonical principal polarization.
On the ordinary Hilbert--Blumenthal moduli locus of these objects the
following assertions hold.

\begin{enumerate}
\item For its geometric generic curve $C$, the Raynaud theta support on
$J=\operatorname{Pic}^0(C^{(1)})$ contains no nonzero point of odd order.
\item For every finite set $S$ of odd primes there is an ordinary such
curve over a finite field $\mathbf F_Q$ with
\[
 \Theta(\overline{\mathbf F}_Q)\cap
 \bigl(J(\overline{\mathbf F}_Q)\{S\}\setminus\{0\}\bigr)
                         =\varnothing.                       \tag{RM.1}
\]
The chosen $\mathcal O$ action and the curve are defined over the same
field $\mathbf F_Q$.  The notation $J\{S\}$ permits all powers and all
products of primes in $S$.
\end{enumerate}

The finite-field examples can be chosen in every prescribed nonempty
open of the ordinary locus with its specified real-multiplication
embedding and principal polarization.  Thus they are Zariski dense
there.  In particular neither $\{3,5\}$ nor any other fixed finite odd
prime support becomes a universal theta-witness supply after the
ordinary real-multiplication reduction of A33.  These assertions are
unconditional; no bound for $Q$ is claimed.

\begin{proof}
\emph{1. The precise moduli input, including characteristic two.}
For an odd integer $h\geq3$, let $H_h$ be the characteristic-two
geometric fibre of the Hilbert--Blumenthal scheme
$M_{(\mathcal O,\mathcal O^+),h}$, with a fixed primitive $h$th root of
unity and full $\mathcal O$-linear symplectic level structure.
Here $\mathcal O^+$ denotes the totally positive elements.
The polarization type $L=\mathcal O$ is exactly an
$\mathcal O$-linear principal polarization: see Deligne--Pappas,
Section~2.6, printed p.~65 \cite{DelignePappas1994}.
Their Theorem~2.2, printed p.~64, gives smoothness since
$2\nmid h\operatorname{disc}(F)=5h$.
Yu's Theorem~1.1, explicitly attributed to Rapoport and
Deligne--Pappas, gives geometric irreducibility for every full level
prime to the residue characteristic \cite[Theorem~1.1, author
manuscript p.~2]{Yu2009HB}.  Thus $H_h$ is smooth and geometrically
irreducible.  The hypotheses require that the level be odd;
they do not require that it be prime to five.

The genus-two Torelli map is an immersion over $\operatorname{Spec}
\mathbf Z$, including characteristic two, by Landesman,
Theorem~1.2 \cite{Landesman2021Torelli}.  In characteristic two its image
is the complement of the product locus $\mathcal A_{1,1}$; the exact
statement is Cl\'ery--van der Geer, Section~4, printed p.~5211
\cite{CleryVanderGeer2022}.  Consequently its pullback to the
smooth, reduced scheme $H_h$ is the open genus-two Jacobian locus,
and this open carries a universal genus-two curve.  More explicitly,
an immersion with this open underlying image becomes the open
subscheme after this base change: its defining ideal has radical
zero on a reduced scheme and is therefore zero.
All genus-two curves are hyperelliptic, so Landesman's
$\mathcal H_2$ is the genus-two curve stack here.

Intersect with the ordinary open and denote the resulting open by
$U_h$.  It is nonempty: the A13 curve, with the integral
Rosati-invariant action proved in the supplement to A33, gives a
geometric point, and a full level-$h$ structure can be chosen over
the algebraic closure.  There is therefore a universal smooth proper
curve $\mathcal C\to U_h$ with its principally polarized Jacobian
$\mathcal A$ and its specified $\mathcal O$ action.  This construction
produces the curve before any finite-field Frobenius specialization;
no extension of the eventual specialization field is needed to
recover a curve from its Jacobian.

\emph{2. Full joint finite-level monodromy and the prime five.}
Fix an odd integer $n>1$ and choose the auxiliary level $h\geq3$
prime to $n$.  Work first over $\overline{\mathbf F}_2$, with compatible
choices of roots of unity.  The forgetful map
\[
                      H_{hn}\longrightarrow H_h             \tag{RM.2}
\]
is the torsor of full $\mathcal O$-linear symplectic level-$n$
structures.  Its group is
\[
                         G_n=\operatorname{SL}_2(\mathcal O/n\mathcal O).
                                                               \tag{RM.3}
\]
Here is the pairing convention, which matters when $5\mid n$.
Let $\mathfrak d^{-1}$ be the inverse different of $F/\mathbf Q$.
The principal polarization determines a perfect alternating
$\mathcal O/n\mathcal O$-bilinear pairing with values in
$\mathfrak d^{-1}\otimes\mu_n$; see Deligne--Pappas,
Sections~2.11--2.13, printed pp.~66--67
\cite{DelignePappas1994}.  After choosing a generator of this
rank-one module, the form is the standard determinant form on a
free rank-two module.  Its $\mathcal O$-linear automorphism group
is (RM.3).  Equivalently, since
$\mathfrak d^{-1}=\frac1{\sqrt5}\mathcal O$, the underlying
$\mathbf Z/n$ pairing is obtained from
\[
 (x,y)\longmapsto
   \operatorname{Tr}_{F/\mathbf Q}
       \bigl((x_1y_2-x_2y_1)/\sqrt5\bigr)\pmod n.             \tag{RM.4}
\]
The inverse-different trace duality makes this form perfect even
at five.  Using the trace pairing on $\mathcal O$ without the
inverse different would not justify that assertion.

Both schemes in (RM.2) are geometrically irreducible by the cited
theorem.  The inverse image of $U_h$, or of any nonempty open in
$U_h$, is a nonempty open of $H_{hn}$ and hence is irreducible.
Therefore the geometric monodromy of $\mathcal A[n]$ on every
such open is all of $G_n$.  This is a joint statement for the
entire integer $n$, so no independence assertion between prime
levels is being assumed.

Put $\mathcal X=\mathcal C^{(1)}$ and
$\mathcal J=\operatorname{Pic}^0(\mathcal X/U_h)$.
Relative Frobenius identifies the odd-torsion local systems of
$\mathcal A$ and $\mathcal J$ by an $\mathcal O$-linear
symplectic similitude of multiplier two.  Conjugation by this
similitude preserves the determinant-one group.  Thus the same
full $G_n$ monodromy holds for $\mathcal J[n]$.

\emph{3. Every nonzero local orbit contains an odd translation cycle.}
Write the finite ring as a product of local principal ideal rings
\[
 \mathcal O/n\mathcal O=\prod_{\mathfrak p\mid n}
                 R_{\mathfrak p},\qquad
 R_{\mathfrak p}=\mathcal O/\mathfrak p^{a_{\mathfrak p}}.
                                                               \tag{RM.5}
\]
Let $z$ be a nonzero element of $(\mathcal O/n\mathcal O)^2$.
Choose a factor in which its component is nonzero, write that
factor as $R=\mathcal O/\mathfrak p^a$, and let $\ell$ be its
odd residue characteristic.  Choose a uniformizer $\varpi$ of
the local ring.  A matrix in $\operatorname{SL}_2(R)$ sends
this component to $\varpi^v e_1$ for some $0\leq v<a$:
divide both coordinates by their largest common power of
$\varpi$, complete the resulting unimodular column to a
determinant-one basis, and invert that basis matrix.

After this change of basis, let $T$ have component
$\varpi^{a-1}e_2$ in the selected factor and zero in every
other factor.  It is nonzero and has additive order $\ell$.
For $j\in\mathbf Z/\ell$, the lower unipotent matrix with
lower-left entry $j\varpi^{a-1-v}$ sends the selected
component of $z$ to
\[
                    \varpi^v e_1+j\varpi^{a-1}e_2.             \tag{RM.6}
\]
Act as the identity in all remaining factors.  It follows that
the $G_n$ orbit of $z$ contains the entire cycle
$z+\langle T\rangle$, consisting of $\ell\geq3$ distinct points.
This proof covers split and inert primes.  It also covers five:
there $\mathcal O/5^r\mathcal O$ is a local ring of length
$2r$ in its uniformizer, and the same argument applies without
any assumption that the ring be reduced.

\emph{4. Generic theta avoidance.}
Let $K$ be the function field of $U_h$ over
$\overline{\mathbf F}_2$ and let $\Theta$ be the Raynaud theta
support of the generic curve.  It is defined over $K$.
If it contained a nonzero point $z\in J[n](\overline K)$,
Galois invariance and Step~2 would put the entire $G_n$ orbit
of $z$ in $\Theta$.  Step~3 then supplies a nonzero $T$ of
odd prime order $\ell$ such that
\[
                   z+\langle T\rangle
                         \subset\Theta\cap(\Theta+T).         \tag{RM.7}
\]
For a genus-two translated Abel curve the intersection with
a distinct nonzero translate has at most two geometric points,
as proved in A19.  Equation (RM.7) contradicts $\ell\geq3$.
This proves absence of all nonzero $n$-torsion at the geometric
generic point for each fixed odd $n$.

The auxiliary level can be chosen prime to each integer under
consideration.  Forgetting this finite auxiliary level does not
change the geometric generic curve or whether its theta contains
a geometric torsion point.  Thus the conclusion holds for every
odd $n$ on the geometric generic curve of the original
real-multiplication locus.  It is not an assertion that a
single open set excludes infinitely many different levels.

\emph{5. A finite avoidance open and an exact Frobenius class.}
Let $S$ be nonempty, put $N=\prod_{\ell\in S}\ell$, and choose
$h$ prime to $N$.  Begin in the inverse image of any prescribed
nonempty ordinary moduli open.  On its universal curve the
relative sequence
\[
 0\longrightarrow\mathcal O_{\mathcal X}
 \longrightarrow F_*\mathcal O_{\mathcal C}
 \longrightarrow\mathcal B\longrightarrow0                 \tag{RM.8}
\]
has a line-bundle quotient of fibre degree one.  As in A22,
the image of the proper relative Abel map
$R\mapsto\mathcal O(R)\otimes\mathcal B^{-1}$ is precisely
the relative Raynaud theta support.  Its intersection with
the finite etale scheme $\mathcal J[N]\setminus\{0\}$ has
closed image on the base.  Step~4 shows that this image
misses the generic point.  Remove it and choose a nonempty
smooth affine open $U$.  Every geometric fibre over $U$ has
\[
                         \Theta_t\cap
                    (\mathcal J_t[N]\setminus\{0\})
                              =\varnothing.                  \tag{RM.9}
\]

Descend the family and this open to a finite field
$k_0=\mathbf F_{q_0}$ large enough that all the data are
defined, $U$ is geometrically integral, and
$q_0\equiv1\pmod{hN^2}$.  The full
$\mathcal O$-linear symplectic level-$N^2$ torsor
\[
                      V\longrightarrow U,
 \qquad G=\operatorname{SL}_2(\mathcal O/N^2\mathcal O),         \tag{RM.10}
\]
is finite etale and geometrically integral by Step~2.
For every $\ell\in S$, set
\[
 u_\ell=1+\ell\pmod{\ell^2},\qquad
 \gamma_\ell=
   \operatorname{diag}(u_\ell,u_\ell^{-1})
       \in\operatorname{SL}_2(\mathcal O/\ell^2\mathcal O).
                                                               \tag{RM.11}
\]
These are $\mathcal O$-linear matrices, including at five.
Each has order $\ell$; by the Chinese remainder theorem
$\gamma=(\gamma_\ell)_{\ell\in S}\in G$ has order $N$.
On the underlying rank-four $\mathbf Z/\ell^2$ module,
\[
 \gamma_\ell\equiv I\pmod\ell,
 \qquad (\gamma_\ell-I)/\ell\pmod\ell
             \text{ is invertible}.                          \tag{RM.12}
\]
Indeed on the rank-two $\mathcal O/\ell$ module the latter
map is $\operatorname{diag}(1,-1)$.  Both entries remain
units when $\mathcal O/5$ is nonreduced.

For completeness, the finite-field realization is exactly the
affine twisting argument of A22, with the group in (RM.10).
Let $\mathcal F$ be arithmetic $q_0$-power Frobenius on $V$.
The action $v\mapsto\mathcal F(v)\gamma^{-1}$ defines a
finite descent datum over $\mathbf F_{q_0^N}/k_0$.
Since $V$ is affine it is effective by the Stacks Project,
Lemma~35.37.1, Tag~\texttt{0245} \cite{Stacks0245}.
The resulting twist is geometrically integral.  Lang--Weil,
in the form of Poonen, Theorem~7.7.1(ii)--(iii), printed
p.~222 \cite{Poonen2017RPV}, supplies a rational point over
$\mathbf F_{q_0^m}$ for every sufficiently large $m$.
Choose $m\equiv1\pmod N$.  Its underlying frame satisfies
\[
              \mathcal F^m(v)=v\gamma^m=v\gamma.               \tag{RM.13}
\]
It lies over $t\in U(\mathbf F_Q)$, $Q=q_0^m$, and the
universal family gives the required ordinary curve and
its specified real multiplication over this exact field.

\emph{6. From a finite level to every exponent and product.}
Let $J$ be the Frobenius-twisted Jacobian of this fibre and
$\Phi$ its $Q$-power Frobenius.  Lifting the torsion frame
to its $\ell$-adic Tate module, (RM.12)--(RM.13) give
\[
             \Phi-I=\ell A_\ell,
             \qquad A_\ell\in\operatorname{GL}_4(\mathbf Z_\ell).
                                                               \tag{RM.14}
\]
Hence, on $(\mathbf Q_\ell/\mathbf Z_\ell)^4$,
$\ker(\Phi-I)=\ker[\ell]$, and therefore
\[
 J(\mathbf F_Q)[\ell^\infty]=J[\ell],\qquad
                       J(\mathbf F_Q)\{S\}=J[N].              \tag{RM.15}
\]
In particular Frobenius fixes every $J[\ell]$ for
$\ell\in S$.  Finding~A19 with $d=1$ now puts every
geometric $S$-primary theta point in this rational group.
Equation (RM.9) excludes its nonzero elements and proves
(RM.1).  If $S$ is empty the assertion is immediate.

The construction began in an arbitrary prescribed nonempty
open, and all subsequent operations preserve a nonempty
open before specialization.  This proves the density
assertion as well.
\end{proof}

\paragraph{Scope of the obstruction.}
A37 rules out a uniform finite odd prime support even inside
the ordinary Rosati-$\mathcal O_5$ locus singled out by A33.
It does not exclude theta witnesses whose prime factors
depend on the individual curve.  It also makes no claim
that the particular explicit curves in A13 or A36 avoid
an arbitrary prescribed support; those curves provide
nonemptiness of the family, whereas the examples in
(RM.1) are obtained by the stated specialization.
